\section{Functions and Roots}

In numerical analysis, the problem of solving a nonlinear equation is commonly expressed as finding a root of a function. A root of a function \( f(x) \) is a value \( x^* \in \mathbb{R} \) such that

\begin{equation}
f(x^*) = 0.
\end{equation}

Depending on the structure of the function, there may exist a single root, multiple roots, or no real roots. In many practical applications, the function \( f \) is continuous, and root-finding methods often rely on this property \cite{burden, atkinson}.

\section{Fixed-Point Formulation}

Many numerical methods are based on transforming the root-finding problem into a fixed-point problem. Instead of solving \( f(x) = 0 \), the equation is rewritten in the form

\begin{equation}
x = g(x),
\end{equation}

where \( g \) is a suitably defined function.

A point \( x^* \) is called a fixed point of \( g \) if

\begin{equation}
x^* = g(x^*).
\end{equation}

If \( g \) satisfies certain conditions, iterative methods of the form

\begin{equation}
x_{k+1} = g(x_k)
\end{equation}

can be used to approximate the solution \cite{burden}.

\section{Convergence of Iterative Methods}

An iterative method generates a sequence \( \{x_k\} \) that is expected to approach the true solution \( x^* \). The method is said to converge if

\begin{equation}
\lim_{k \to \infty} x_k = x^*.
\end{equation}

The convergence of an iterative method depends on several factors, including the choice of the initial guess and the properties of the function.

A common condition for convergence in fixed-point iteration is that the function \( g \) is a contraction, meaning that there exists a constant \( L < 1 \) such that

\begin{equation}
|g(x) - g(y)| \leq L |x - y|
\end{equation}

for all \( x, y \) in a neighborhood of the fixed point \cite{burden, cs357}.

\section{Order of Convergence}

The speed at which an iterative method converges is described by its order of convergence. Let \( e_k = |x_k - x^*| \) denote the error at iteration \( k \). The method is said to have order \( p \) if there exists a constant \( C > 0 \) such that

\begin{equation}
e_{k+1} \approx C e_k^p
\end{equation}

for sufficiently large \( k \).

For example:
\begin{itemize}
    \item \( p = 1 \): linear convergence
    \item \( p = 2 \): quadratic convergence
\end{itemize}

Higher-order methods generally converge faster, but they may require more computational effort per iteration \cite{burden}.

\section{Norms and Vector Functions}

For systems of nonlinear equations, it is necessary to measure the size of vectors. A commonly used norm is the Euclidean norm, defined as

\begin{equation}
\|\mathbf{x}\| = \sqrt{x_1^2 + x_2^2 + \cdots + x_n^2}.
\end{equation}

Norms are used to quantify errors and convergence in higher-dimensional problems \cite{atkinson}.

\section{Jacobian Matrix}

For a system of nonlinear equations \( F(\mathbf{x}) = \mathbf{0} \), the Jacobian matrix plays a central role in many numerical methods.

The Jacobian matrix is defined as

\begin{equation}
J_F(\mathbf{x}) =
\begin{bmatrix}
\frac{\partial f_1}{\partial x_1} & \cdots & \frac{\partial f_1}{\partial x_n} \\
\vdots & \ddots & \vdots \\
\frac{\partial f_n}{\partial x_1} & \cdots & \frac{\partial f_n}{\partial x_n}
\end{bmatrix}.
\end{equation}

It generalizes the concept of the derivative to multivariable functions and is essential in methods such as Newton’s method for systems \cite{kelley}.

\section{Stopping Criteria}

In practical computations, iterative methods cannot run indefinitely. Therefore, a stopping criterion is required.

Common stopping criteria include:
\begin{itemize}
    \item \( |x_{k+1} - x_k| < \varepsilon \)
    \item \( |f(x_k)| < \varepsilon \)
    \item maximum number of iterations reached
\end{itemize}

where \( \varepsilon \) is a predefined tolerance \cite{burden, cs357}.